\section{Introduction}\label{sec:intro}

\IEEEPARstart{L}{arge} language model (LLM) inference has shifted cloud datacenter energy demand from offline batch training toward continuous, latency-critical serving. Survey estimates indicate that U.S.\ datacenter electricity consumption could climb from \SI{176}{\tera\watt\hour} in 2023 toward \SI{580}{\tera\watt\hour} by 2028~\cite{yang2025container}. Because interactive LLM serving runs across distributed cloud clusters, its net carbon footprint depends on compute power draw, temporal and regional electricity grid carbon intensity ($\text{gCO}_2\text{eq/kWh}$), and hardware manufacturing emissions amortized over server and accelerator lifetimes~\cite{gupta2021chasing,iso2024sci}.

Two distinct research communities have tackled pieces of this problem over the past five years, yet their mechanisms remain only partially connected. First, systems researchers have shown that autoregressive LLM generation divides into a compute-bound \emph{prefill} phase (governing time-to-first-token, TTFT) and a memory-bandwidth-bound \emph{decode} phase (governing time-per-output-token, TPOT)~\cite{patel2024splitwise,zhong2024distserve,agrawal2024sarathi}. Disaggregating prefill and decode across separate GPU pools eliminates phase interference and enables power capping or lower-tier decode GPUs~\cite{patel2024splitwise,stojkovic2025dynamollm,basit2026dualscale,yan2026pdprov}. Second, sustainable cloud researchers have built carbon-aware spatial/temporal routers and serverless keep-alive controllers that follow low-carbon regions or clean grid hours~\cite{wiesner2021letswait,hanafy2023carbonscaler,sukprasert2024limitations,gsteiger2024caribou,sun2026lacere}.

A PRISMA-2020-guided literature review of 65 retained works (51 empirical/modeling studies, 11 surveys or agendas, 2 standards, and 1 production trace release) clarifies both recent 2024--2026 progress and the remaining architectural gap. On single clusters, GreenLLM~\cite{shi2024greenllm} pairs heterogeneous P/D disaggregation with grid CI and GPU lifetime reuse, while SPROUT~\cite{li2024sprout} and GreenCache~\cite{tian2026greencache} model operational plus hardware/SSD embodied carbon. Across multiple regions, recent geo-distributed LLM routers~\cite{topcu2026tension,moore2025slit,zhang2026context,bernhard2026cleanest,he2025freesh,more2026cedar,cheng2026greenllm} cut operational emissions over colocated pools. Yet across all 65 works, 50/65 (76.9\%) omit manufacturing emissions, 0/11 implemented cross-region carbon schedulers evaluate separated P/D pools with wide-area KV transfer gating (FREESH~\cite{he2025freesh} uses P/D notation in Eq.~1c while evaluating colocated pools, and CarbonRouter~\cite{zhang2026context} defers WAN KV migration), InFactPlanner~\cite{tsiopani2026infactplanner} abstracts away batching, KV state, and P/D separation, and Sisodia's \texttt{llm-d} agenda~\cite{sisodia2026sicp} has zero experiments. Thus, no implemented system jointly integrates P/D disaggregation, cross-region carbon routing with WAN KV gating, host plus GPU embodied accounting, and serverless keep-alive across 10.3--84.0~s 70B loader tiers~\cite{fu2024serverlessllm}.

To evaluate this joint intersection under controlled, reproducible conditions, I conduct a PRISMA-2020-guided review paired with a trace-driven discrete-event simulation study structured around four research questions:
\begin{itemize}
    \item \textbf{RQ1 (Taxonomy and Joint-Design Gaps):} Where do existing carbon-aware cloud schedulers, geo-distributed LLM routers, serverless keep-alive controllers, and phase-disaggregated serving engines intersect, and what remains unimplemented at their boundary?
    \item \textbf{RQ2 (Disaggregated Carbon--SLO Frontier):} Under real UK/Japan grid telemetry and production Azure LLM traces, when does hybrid carbon-and-queue-aware phase routing outperform latency-only colocation and phase-agnostic or unconditional spatial routing under TTFT/TPOT SLOs?
    \item \textbf{RQ3 (Serverless Keep-Alive Economics):} Across carbon-intensive and wind-clean grids, how does hazard-driven dynamic keep-alive trade operational carbon against cold-start SLO violations as loader latency spans fast NVMe streaming to standard deserialization?
    \item \textbf{RQ4 (Embodied Scale-Out Paradox):} When does adding spare GPU nodes in clean regions invert net lifecycle Software Carbon Intensity ($\text{SCI}_{\text{tot}}$), and how do decoupled host/GPU lifetimes (9~yr/3~yr vs.\ 4~yr) alter provisioning decisions?
\end{itemize}

This paper makes three contributions:
\begin{enumerate}
    \item \textbf{PRISMA-2020-Guided Literature Synthesis:} Screening 102 records under explicit single-author plus AI-assisted review disclosure yields a verified 65-work corpus (60 full-text, 5 abstract-only), a dual-axis taxonomy (\texttt{L1.A--L1.F}), a direct feature comparison against the closest 2024--2026 systems, and seven open challenges (\texttt{OC1--OC7}).
    \item \textbf{Benchmark-Validated Carbon-Extended Simulator:} I formalize an ISO/IEC 21031:2024~\cite{iso2024sci} lifecycle model inside SimPy 4.1.2 integrating Splitwise power/step curves~\cite{patel2024splitwise}, GreenSKU/LLMCarbon host-plus-GPU embodied parameters~\cite{li2025ecoserve,faiz2024llmcarbon}, WAN KV transfer-vs-recompute gating, and GapHazard serverless control~\cite{sun2026lacere,fu2024serverlessllm}, separated into 19 calibration anchors and 5 held-out cross-load validation checks against published baselines.
    \item \textbf{Multi-Seed Trace-Driven Empirical Evaluation:} Across 10 paired seeds using CC-BY-4.0 Azure traces~\cite{azure2024trace}, UK NESO, and Japan TSO telemetry (with 24h trace-tiling sensitivity ablation), hybrid routing (\texttt{P4}), hazard keep-alive (\texttt{S3}), and elastic over-provisioning (\texttt{H4}) quantify the joint carbon--SLO--embodied trade-offs detailed in Section~\ref{sec:results}.
\end{enumerate}

Section~\ref{sec:background} reviews preliminaries; Sections~\ref{sec:methodology} and~\ref{sec:taxonomy} present the PRISMA review, taxonomy (\texttt{RQ1}), and related work; Sections~\ref{sec:sim_design} and~\ref{sec:results} detail the simulator and multi-seed findings (\texttt{RQ1--RQ4}); Section~\ref{sec:challenges} discusses open challenges and limitations; and Section~\ref{sec:conclusion} concludes with disclosures.
